% p212_1; source PDF page 212; manually verified editable table.
\begin{tikzpicture}[x=1cm,y=-1cm,inner sep=0pt,outer sep=0pt]
\path[use as bounding box] (0,0) rectangle (16.000,4.350);
\draw[line width=.4pt] (0.000,0.000) rectangle (1.800,0.850);
\node[align=center,text width=1.540cm,font={\linespread{1}\fontsize{10}{13.50}\selectfont\bfseries}] at (0.900,0.425) {追诉期限};
\draw[line width=.4pt] (1.800,0.000) rectangle (4.400,0.850);
\node[align=center,text width=2.340cm,font={\linespread{1}\fontsize{10}{13.50}\selectfont\bfseries}] at (3.100,0.425) {计算标准};
\draw[line width=.4pt] (4.400,0.000) rectangle (16.000,0.850);
\node[align=center,text width=11.340cm,font={\linespread{1}\fontsize{10}{13.50}\selectfont\bfseries}] at (10.200,0.425) {要点（2015、2018、2023年试题）};
\draw[line width=.4pt] (0.000,0.850) rectangle (1.800,3.450);
\node[align=center,text width=1.540cm,font={\linespread{1}\fontsize{10}{13.50}\selectfont}] at (0.900,2.150) {起算日};
\draw[line width=.4pt] (1.800,0.850) rectangle (4.400,3.450);
\node[align=center,text width=2.340cm,font={\linespread{1}\fontsize{10}{13.50}\selectfont}] at (3.100,2.150) {犯罪成立之日};
\draw[line width=.4pt] (4.400,0.850) rectangle (16.000,3.450);
\node[align=left,text width=11.340cm,font={\linespread{1}\fontsize{10}{13.50}\selectfont}] at (10.200,2.150) {1. 故意犯，实施行为，产生危险，犯罪就成立。造成实害结果，犯罪既遂。\\2. 过失犯，造成实害结果，犯罪成立。\\例外：连续犯、继续犯。犯罪行为有连续或者继续状态的，从犯罪行为终了之日起计算（2025年试题，非法持有枪支罪）};
\draw[line width=.4pt] (0.000,3.450) rectangle (1.800,4.350);
\node[align=center,text width=1.540cm,font={\linespread{1}\fontsize{10}{13.50}\selectfont}] at (0.900,3.900) {截止日};
\draw[line width=.4pt] (1.800,3.450) rectangle (4.400,4.350);
\node[align=center,text width=2.340cm,font={\linespread{1}\fontsize{10}{13.50}\selectfont}] at (3.100,3.900) {审判之日};
\draw[line width=.4pt] (4.400,3.450) rectangle (16.000,4.350);
\node[align=left,text width=11.340cm,font={\linespread{1}\fontsize{10}{13.50}\selectfont}] at (10.200,3.900) {只要在审判之日还没有超过追诉期限，就能追诉};
\end{tikzpicture}
